% Fichier complété par tools/complete_missing_corrections.py.
% Source : ELEC/ELEC-01/538_CircuitElec/538_CircuitElec.tex
% Statut : rédigé (2 question(s)).
\begin{corrigebox}[Corrigé rédigé]
\CorrigeQuestion{1}
On réduit le réseau depuis l'extrémité opposée aux bornes. Deux
            résistances en série donnent $R_s=R_a+R_b$ et deux résistances en parallèle
            $R_p=R_aR_b/(R_a+R_b)$. Pour le réseau en échelle, on définit récursivement
            $Z_{n}=R_{v,n}$ puis
            \[Z_k=R_{v,k}\parallel(R_{h,k+1}+Z_{k+1}),\qquad
              R_{eq}=R_{h,1}+Z_1.\]
            Pour le pont symétrique, les nœuds équipotentiels peuvent être fusionnés avant
            la réduction.
\CorrigeQuestion{2}
Le courant fourni par la source est $I=E/R_{eq}$. On
            remonte ensuite le réseau : la tension d'un bloc parallèle est commune à ses
            branches et les courants se partagent selon
            $I_a=U/R_a$, $I_b=U/R_b$. À chaque étage, la loi des nœuds vérifie
            $I_{amont}=I_{vertical}+I_{aval}$.
\end{corrigebox}
